\documentclass[10pt]{beamer}
\usetheme{metropolis}
\usepackage{graphicx}
\usepackage{listings}
\usepackage{booktabs}
\usepackage{xcolor}
\usepackage{hyperref}

\definecolor{codebg}{HTML}{F5F5F5}
\definecolor{codegreen}{HTML}{2E7D32}
\definecolor{codegray}{HTML}{757575}
\definecolor{codeblue}{HTML}{1565C0}

\lstdefinestyle{pythonstyle}{
  language=Python, backgroundcolor=\color{codebg},
  basicstyle=\ttfamily\scriptsize,
  keywordstyle=\color{codeblue}\bfseries,
  stringstyle=\color{codegreen},
  commentstyle=\color{codegray}\itshape,
  breaklines=true, frame=single,
  rulecolor=\color{codegray!30}, showstringspaces=false, tabsize=4,
  morekeywords={import,from,as,pd,plt,sns,np,fig,ax,subplots,savefig,
    scatter,hist,barh,boxplot,set_xlabel,set_ylabel,set_title,
    spines,set_visible,tick_params,tight_layout,DataFrame,
    read_csv,head,info,describe,dtypes,shape,isna,sum,
    value_counts,groupby,agg,sort_values,reset_index}
}
\lstdefinestyle{bashstyle}{
  language=bash, backgroundcolor=\color{codebg},
  basicstyle=\ttfamily\scriptsize,
  keywordstyle=\color{codeblue}\bfseries,
  commentstyle=\color{codegray}\itshape,
  breaklines=true, frame=single,
  rulecolor=\color{codegray!30}, showstringspaces=false,
}

\title{Python Environment Setup \& Matplotlib Basics}
\subtitle{Workshop 1 --- Module 9}
\author{Prof.Asoc.\ Endri Raco}
\institute{Data Visualization for Data Scientists \\ UNYT Tirana}
\date{2025/26}

\begin{document}

\begin{frame}
  \titlepage
\end{frame}

% ================================================================
\begin{frame}{Module Roadmap}
  \begin{enumerate}
    \item Installing Python, JupyterLab, and core packages
    \item The Python data science stack: NumPy, pandas, matplotlib, seaborn
    \item pandas data types and the DataFrame inspection workflow
    \item Matplotlib architecture: Figure $\rightarrow$ Axes $\rightarrow$ Artists
    \item The four fundamental chart types: scatter, histogram, bar, boxplot
    \item pyplot vs Object-Oriented API
  \end{enumerate}
  \begin{exampleblock}{Learning Outcome}
    You will have a functioning Python environment and be able to produce
    all four fundamental chart types using matplotlib, preparing the
    ground for seaborn/plotnine in Workshop 2.
  \end{exampleblock}
\end{frame}

% ================================================================
\section{Python \& JupyterLab Setup}
% ================================================================

\begin{frame}[fragile]{Installation Checklist}
  \begin{columns}[T]
    \begin{column}{0.48\textwidth}
      \textbf{Step 1: Install Python}
\begin{lstlisting}[style=bashstyle]
# Option A: Miniconda (recommended)
# https://docs.conda.io/en/latest/
conda create -n dataviz python=3.11
conda activate dataviz

# Option B: pip + venv
python -m venv dataviz_env
source dataviz_env/bin/activate
\end{lstlisting}
      \vspace{4pt}
      \textbf{Step 2: Install JupyterLab}
\begin{lstlisting}[style=bashstyle]
pip install jupyterlab
jupyter lab  # opens browser
\end{lstlisting}
    \end{column}
    \begin{column}{0.48\textwidth}
      \textbf{Step 3: Install core packages}
\begin{lstlisting}[style=bashstyle]
pip install \
  numpy pandas \
  matplotlib seaborn \
  plotnine scipy \
  openpyxl  # for Excel files
\end{lstlisting}
      \vspace{4pt}
      \textbf{Verify installation:}
\begin{lstlisting}[style=pythonstyle]
import numpy as np
import pandas as pd
import matplotlib.pyplot as plt
import seaborn as sns

print(f"numpy:  {np.__version__}")
print(f"pandas: {pd.__version__}")
print(f"mpl:    {plt.matplotlib
                    .__version__}")
print(f"sns:    {sns.__version__}")
\end{lstlisting}
    \end{column}
  \end{columns}
\end{frame}

% ================================================================
\begin{frame}{The JupyterLab IDE}
  \begin{center}
    \includegraphics[width=0.88\textwidth]{figures_m09/jupyter_layout}
  \end{center}
\end{frame}

% ================================================================
\begin{frame}[fragile]{JupyterLab Keyboard Shortcuts}
  \centering
  \begin{tabular}{lll}
    \toprule
    \textbf{Action} & \textbf{Command Mode} & \textbf{Edit Mode} \\
    \midrule
    Run cell + next           & Shift+Enter & Shift+Enter \\
    Run cell (stay)           & Ctrl+Enter  & Ctrl+Enter \\
    Insert cell above         & A           & --- \\
    Insert cell below         & B           & --- \\
    Delete cell               & D, D        & --- \\
    Switch to Markdown        & M           & --- \\
    Switch to Code            & Y           & --- \\
    Toggle line numbers       & L           & --- \\
    \bottomrule
  \end{tabular}

  \vspace{6pt}
  \begin{alertblock}{Pro-Tip}
    Press \textbf{Escape} to enter Command Mode (blue sidebar),
    \textbf{Enter} to enter Edit Mode (green sidebar). Most navigation
    happens in Command Mode; typing happens in Edit Mode.
  \end{alertblock}
\end{frame}

% ================================================================
\section{The Python Data Science Stack}
% ================================================================

\begin{frame}{Python Ecosystem Overview}
  \begin{center}
    \includegraphics[width=0.88\textwidth]{figures_m09/python_stack}
  \end{center}
\end{frame}

% ================================================================
\begin{frame}{R $\leftrightarrow$ Python Package Equivalents}
  \centering
  \begin{tabular}{lll}
    \toprule
    \textbf{Task} & \textbf{R} & \textbf{Python} \\
    \midrule
    Data frames    & \texttt{tibble / data.frame}   & \texttt{pandas.DataFrame} \\
    CSV import     & \texttt{readr::read\_csv}       & \texttt{pd.read\_csv} \\
    Data wrangling & \texttt{dplyr}                  & \texttt{pandas} \\
    Static plots   & \texttt{ggplot2}                & \texttt{matplotlib / seaborn} \\
    GoG clone      & \texttt{ggplot2} (native)       & \texttt{plotnine} \\
    Interactive    & \texttt{plotly / shiny}          & \texttt{plotly / dash} \\
    Statistics     & built-in + \texttt{broom}       & \texttt{scipy / statsmodels} \\
    ML             & \texttt{tidymodels}             & \texttt{scikit-learn} \\
    \bottomrule
  \end{tabular}

  \vspace{6pt}
  \begin{exampleblock}{Data Insight}
    The ecosystems are deliberately parallel. Skills in one transfer to
    the other with mostly syntactic differences.
  \end{exampleblock}
\end{frame}

% ================================================================
\section{pandas Data Types}
% ================================================================

\begin{frame}{Data Types for Visualization}
  \begin{center}
    \includegraphics[width=0.82\textwidth]{figures_m09/pandas_types}
  \end{center}
  \begin{alertblock}{Pro-Tip}
    Use \texttt{df.dtypes} to check all column types and \texttt{df.info()}
    for a compact summary including non-null counts. Type mismatches
    (e.g., numbers parsed as \texttt{object}) are the number-one cause
    of plotting failures.
  \end{alertblock}
\end{frame}

% ================================================================
\begin{frame}[fragile]{Loading and Inspecting Data}
\begin{lstlisting}[style=pythonstyle]
import pandas as pd

# Read CSV
apps = pd.read_csv("googleplaystore.csv")

# Quick inspection (parallel to R's glimpse/str/summary)
print(apps.shape)         # (rows, cols)  — like dim()
print(apps.dtypes)        # column types  — like str()
print(apps.head())        # first 5 rows  — like head()
print(apps.describe())    # summary stats — like summary()
print(apps.info())        # types + nulls — like glimpse()

# Missing values
print(apps.isna().sum())  # per-column NA count

# Useful for viz planning:
# - How many rows? (affects point density / alpha)
# - dtypes? (object = probably categorical or needs conversion)
# - Nulls? (need dropna() or fillna() before plotting)
\end{lstlisting}
\end{frame}

% ================================================================
\begin{frame}[fragile]{Type Conversion Patterns}
\begin{lstlisting}[style=pythonstyle]
# String to numeric (like as.numeric() in R)
apps["Reviews"] = pd.to_numeric(apps["Reviews"], errors="coerce")

# String to datetime (like as.Date() in R)
apps["Last Updated"] = pd.to_datetime(apps["Last Updated"],
                                       format="%B %d, %Y",
                                       errors="coerce")

# String to categorical (like factor() in R)
apps["Type"] = pd.Categorical(apps["Type"])

# Ordered category (like ordered factor)
rating_order = ["Everyone", "Teen", "Mature 17+"]
apps["Content Rating"] = pd.Categorical(
    apps["Content Rating"],
    categories=rating_order,
    ordered=True)

# Verify
print(apps.dtypes)
\end{lstlisting}
\end{frame}

% ================================================================
\section{Matplotlib Architecture}
% ================================================================

\begin{frame}{Figure $\rightarrow$ Axes $\rightarrow$ Artists}
  \begin{center}
    \includegraphics[width=0.85\textwidth]{figures_m09/mpl_architecture}
  \end{center}
\end{frame}

% ================================================================
\begin{frame}{pyplot vs Object-Oriented API}
  \begin{center}
    \includegraphics[width=0.92\textwidth]{figures_m09/pyplot_vs_oo}
  \end{center}
  \begin{exampleblock}{Data Insight}
    The \textbf{pyplot} API (\texttt{plt.plot()}) uses an implicit state
    machine modelled after MATLAB. The \textbf{OO} API (\texttt{ax.plot()})
    gives explicit control over every figure element. This course uses
    the OO API exclusively --- it is clearer, more composable, and the
    only viable approach for multi-panel layouts.
  \end{exampleblock}
\end{frame}

% ================================================================
\begin{frame}[fragile]{The OO Pattern: Your Template for Every Plot}
\begin{lstlisting}[style=pythonstyle]
import matplotlib.pyplot as plt

# 1. Create figure and axes (ALWAYS start here)
fig, ax = plt.subplots(figsize=(7, 5))

# 2. Plot data onto the axes
ax.scatter(x, y, s=30, c="#1565C0", alpha=0.6)

# 3. Labels and title
ax.set_xlabel("X Variable", fontsize=10)
ax.set_ylabel("Y Variable", fontsize=10)
ax.set_title("My Chart Title", fontsize=12, fontweight="bold")

# 4. Tufte cleanup
ax.spines["top"].set_visible(False)
ax.spines["right"].set_visible(False)
ax.tick_params(labelsize=8)

# 5. Save (ALWAYS use bbox_inches="tight")
plt.tight_layout()
fig.savefig("chart.pdf", bbox_inches="tight")     # vector
fig.savefig("chart.png", dpi=300, bbox_inches="tight")  # raster
plt.close()  # free memory
\end{lstlisting}
\end{frame}

% ================================================================
\section{Four Fundamental Chart Types}
% ================================================================

\begin{frame}{The Four Charts in matplotlib}
  \begin{center}
    \includegraphics[width=0.88\textwidth]{figures_m09/mpl_subplots}
  \end{center}
  \begin{exampleblock}{Data Insight}
    These four methods --- \texttt{ax.scatter()}, \texttt{ax.hist()},
    \texttt{ax.barh()}, \texttt{ax.boxplot()} --- cover 80\% of
    exploratory needs. Master these before moving to seaborn.
  \end{exampleblock}
\end{frame}

% ================================================================
\begin{frame}[fragile]{Scatterplot: \texttt{ax.scatter()}}
  \begin{columns}[T]
    \begin{column}{0.52\textwidth}
\begin{lstlisting}[style=pythonstyle]
fig, ax = plt.subplots(figsize=(5, 4))

ax.scatter(x, y,
    s=30,             # marker size
    c="#1565C0",      # colour
    alpha=0.6,        # transparency
    edgecolors="white",
    linewidth=0.3)

# Add regression line
m, b = np.polyfit(x, y, 1)
xs = np.linspace(x.min(), x.max(), 50)
ax.plot(xs, m*xs + b,
    color="#E53935", lw=1.5,
    linestyle="--")

ax.set_xlabel("Displacement (L)")
ax.set_ylabel("MPG")
ax.set_title("Scatter + OLS Fit")
ax.spines["top"].set_visible(False)
ax.spines["right"].set_visible(False)
\end{lstlisting}
    \end{column}
    \begin{column}{0.45\textwidth}
      \includegraphics[width=\textwidth]{figures_m09/mpl_scatter}

      \vspace{4pt}
      {\small Key params:
      \texttt{s} = size (area in pts\textsuperscript{2}),
      \texttt{c} = colour (string or array),
      \texttt{alpha} = transparency [0--1],
      \texttt{edgecolors} = marker border.}
    \end{column}
  \end{columns}
\end{frame}

% ================================================================
\begin{frame}[fragile]{Histogram: \texttt{ax.hist()}}
  \begin{columns}[T]
    \begin{column}{0.52\textwidth}
\begin{lstlisting}[style=pythonstyle]
fig, ax = plt.subplots(figsize=(5, 4))

ax.hist(data,
    bins=25,              # bin count
    color="#1565C0",
    edgecolor="white",
    linewidth=0.5,
    alpha=0.8)

# Mean reference line
mean_val = np.mean(data)
ax.axvline(x=mean_val,
    color="#E53935", lw=2,
    linestyle="--")
ax.text(mean_val + 1,
    ax.get_ylim()[1] * 0.9,
    f"Mean: {mean_val:.1f}",
    fontsize=8, color="#E53935",
    fontweight="bold")

ax.set_xlabel("Value")
ax.set_ylabel("Frequency")
\end{lstlisting}
    \end{column}
    \begin{column}{0.45\textwidth}
      \includegraphics[width=\textwidth]{figures_m09/mpl_hist}

      \vspace{4pt}
      {\small Key params:
      \texttt{bins} = int or sequence,
      \texttt{density=True} for PDF,
      \texttt{histtype="step"} for unfilled,
      \texttt{cumulative=True} for CDF.}
    \end{column}
  \end{columns}
\end{frame}

% ================================================================
\begin{frame}[fragile]{Bar Chart: \texttt{ax.barh()}}
  \begin{columns}[T]
    \begin{column}{0.52\textwidth}
\begin{lstlisting}[style=pythonstyle]
fig, ax = plt.subplots(figsize=(5, 4))

# Sort ascending for horizontal
idx = np.argsort(vals)
s_cats = [cats[i] for i in idx]
s_vals = [vals[i] for i in idx]

ax.barh(s_cats, s_vals,
    color="#1565C0",
    height=0.55,       # bar width
    edgecolor="none")

# Direct labels
for i, v in enumerate(s_vals):
    ax.text(v + 0.8, i, str(v),
        va="center", fontsize=9,
        fontweight="bold")

# Tufte cleanup
for sp in ["top","right","left"]:
    ax.spines[sp].set_visible(False)
ax.set_xticks([])
ax.tick_params(left=False)
\end{lstlisting}
    \end{column}
    \begin{column}{0.45\textwidth}
      \includegraphics[width=\textwidth]{figures_m09/mpl_bar}

      \vspace{4pt}
      {\small Use \texttt{ax.barh()} (horizontal)
      by default: labels are readable
      without rotation. \texttt{ax.bar()} for
      vertical only when x-axis is
      temporal or ordinal.}
    \end{column}
  \end{columns}
\end{frame}

% ================================================================
\begin{frame}[fragile]{Boxplot: \texttt{ax.boxplot()}}
  \begin{columns}[T]
    \begin{column}{0.52\textwidth}
\begin{lstlisting}[style=pythonstyle]
fig, ax = plt.subplots(figsize=(5, 4))

bp = ax.boxplot(
    [group_a, group_b, group_c],
    labels=["A", "B", "C"],
    patch_artist=True,  # fill boxes
    widths=0.5,
    notch=True)         # CI on median

# Customise box colours
fills = ["#E3F2FD","#E8F5E9","#FFF3E0"]
edges = ["#1565C0","#2E7D32","#E65100"]
for patch, fc, ec in zip(
    bp["boxes"], fills, edges):
    patch.set_facecolor(fc)
    patch.set_edgecolor(ec)

# Overlay mean points
means = [np.mean(d) for d in data]
ax.scatter(range(1, 4), means,
    c="#E53935", s=50, zorder=5,
    marker="D", edgecolors="white")
\end{lstlisting}
    \end{column}
    \begin{column}{0.45\textwidth}
      \includegraphics[width=\textwidth]{figures_m09/mpl_box}

      \vspace{4pt}
      {\small Key params:
      \texttt{patch\_artist=True} enables
      fill colour;
      \texttt{notch=True} adds median CI;
      \texttt{whis=1.5} (default) controls
      whisker extent;
      use \texttt{vert=False} for horizontal.}
    \end{column}
  \end{columns}
\end{frame}

% ================================================================
\section{Multi-Panel Layouts}
% ================================================================

\begin{frame}[fragile]{\texttt{plt.subplots(rows, cols)}}
\begin{lstlisting}[style=pythonstyle]
# Create 2x2 grid (equivalent to R's par(mfrow = c(2,2)))
fig, axes = plt.subplots(2, 2, figsize=(8, 6))

# Access each axes by index
axes[0, 0].scatter(x, y, s=20, c="#1565C0")
axes[0, 0].set_title("(a) Scatter")

axes[0, 1].hist(data, bins=20, color="#2E7D32")
axes[0, 1].set_title("(b) Histogram")

axes[1, 0].barh(cats, vals, color="#E65100", height=0.5)
axes[1, 0].set_title("(c) Bar Chart")

axes[1, 1].boxplot(groups, patch_artist=True)
axes[1, 1].set_title("(d) Boxplot")

# Tufte cleanup on all panels
for ax in axes.flatten():
    ax.spines["top"].set_visible(False)
    ax.spines["right"].set_visible(False)
    ax.tick_params(labelsize=7)

plt.tight_layout()  # prevent label overlap
fig.savefig("multipanel.pdf", bbox_inches="tight")
\end{lstlisting}
\end{frame}

% ================================================================
\begin{frame}[fragile]{Advanced: \texttt{GridSpec} for Unequal Panels}
\begin{lstlisting}[style=pythonstyle]
from matplotlib.gridspec import GridSpec

fig = plt.figure(figsize=(10, 6))
gs = GridSpec(2, 3, figure=fig,
    hspace=0.4, wspace=0.3,
    height_ratios=[1, 1.5])

# Top row: three equal panels
ax1 = fig.add_subplot(gs[0, 0])  # row 0, col 0
ax2 = fig.add_subplot(gs[0, 1])  # row 0, col 1
ax3 = fig.add_subplot(gs[0, 2])  # row 0, col 2

# Bottom row: one full-width panel
ax4 = fig.add_subplot(gs[1, :])  # row 1, ALL cols

# Shared y-axis: sharey=ax1
ax5 = fig.add_subplot(gs[0, 2], sharey=ax1)
\end{lstlisting}

  \begin{alertblock}{Pro-Tip}
    \texttt{GridSpec} is the Python equivalent of R's \texttt{patchwork}
    package. Use it for any layout that is not a regular grid.
    Slice notation (\texttt{gs[1, :]}) spans multiple columns.
  \end{alertblock}
\end{frame}

% ================================================================
\section{Synthesis}
% ================================================================

\begin{frame}{R vs Python: Side-by-Side Cheat Sheet}
  \centering
  {\small
  \begin{tabular}{lll}
    \toprule
    \textbf{Operation} & \textbf{R (base)} & \textbf{Python (matplotlib)} \\
    \midrule
    Scatter   & \texttt{plot(x, y)}             & \texttt{ax.scatter(x, y)} \\
    Histogram & \texttt{hist(x)}                & \texttt{ax.hist(x)} \\
    Bar chart & \texttt{barplot(table(x))}      & \texttt{ax.barh(cats, vals)} \\
    Boxplot   & \texttt{boxplot(y \textasciitilde{} g)} & \texttt{ax.boxplot([g1, g2])} \\
    Multi-panel & \texttt{par(mfrow=c(2,2))}    & \texttt{plt.subplots(2, 2)} \\
    Save plot & \texttt{png(); ...; dev.off()}  & \texttt{fig.savefig("f.png")} \\
    Reg.\ line & \texttt{abline(lm(y\textasciitilde{}x))} & \texttt{ax.plot(xs, m*xs+b)} \\
    Ref.\ line & \texttt{abline(h=v)}           & \texttt{ax.axhline(y=v)} \\
    \bottomrule
  \end{tabular}}
\end{frame}

% ================================================================
\begin{frame}{Key Takeaways}
  \begin{enumerate}
    \item \textbf{Miniconda + JupyterLab} is the recommended Python
          setup. Learn Shift+Enter and Command/Edit mode switching.
    \item \textbf{pandas} is the Python equivalent of tibble + dplyr.
          Use \texttt{df.dtypes} and \texttt{df.info()} before plotting.
    \item Matplotlib has \textbf{three layers}: Figure (canvas), Axes
          (plot area), Artists (data marks). Always think in this hierarchy.
    \item Use the \textbf{OO API} (\texttt{fig, ax = plt.subplots()})
          exclusively --- it is explicit, composable, and scalable.
    \item Master four methods: \texttt{ax.scatter()}, \texttt{ax.hist()},
          \texttt{ax.barh()}, \texttt{ax.boxplot()}.
    \item \texttt{plt.subplots(r, c)} for regular grids;
          \texttt{GridSpec} for unequal panels.
  \end{enumerate}
\end{frame}

% ================================================================
\begin{frame}{Essential Readings for This Module}
  \begin{itemize}
    \item VanderPlas, J. (2016). \emph{Python Data Science Handbook},
          Chapter 4 (Visualization with Matplotlib).
          \texttt{https://jakevdp.github.io/PythonDataScienceHandbook/}
    \item McKinney, W. (2022). \emph{Python for Data Analysis}, 3rd ed.,
          Chapters 4--5 (NumPy, pandas) and Chapter 9 (Plotting).
    \item Matplotlib official tutorials:
          \texttt{https://matplotlib.org/stable/tutorials/}
  \end{itemize}

  \vspace{6pt}
  \textbf{Next Module}: Comparative Lab --- First Plots in R vs Python (W01-M10)
\end{frame}

\end{document}
